\section{Tricks to Compute Integrals Easily}

\begin{itemize}
    \item \textbf{Cauchy's Theorem:} If $f$ is holomorphic on a \textbf{simply connected} set $D$, then for any closed loop $\gamma$ in $D$:
    \begin{align*}
        \int_{\gamma} f(z) \, dz = 0
    \end{align*}
    
    \item \textbf{Deformation Theorem:} We can replace an integral over a "complicated" curve $\Gamma_0$ by an integral over a "simpler" curve $\Gamma_1$ (like a small circle) as long as $f$ is holomorphic in the region between $\Gamma_0$ and $\Gamma_1$.
    
    \item \textbf{Path Independence:} If $\Gamma_0$ and $\Gamma_1$ share the same endpoints and orientation, and $f$ is holomorphic on the region they enclose, then:
    \begin{align*}
        \int_{\Gamma_0} f(z) \, dz = \int_{\Gamma_1} f(z) \, dz
    \end{align*}
    
    \item \textbf{Generalized Cauchy Integral Formula:} Useful for computing integrals of the form $\frac{f(z)}{(z-z_0)^{n+1}}$. If $f$ is holomorphic on and inside $\gamma$:
    \begin{align*}
        \int_\gamma \frac{f(z)}{(z-z_0)^{n+1}} \, dz = \frac{2 \pi i}{n!} f^{(n)}(z_0)
    \end{align*}
    \textit{Example:} To calculate $\int_\gamma \frac{e^z}{z^3} dz$ around $0$, use $n=2$ and $f(z)=e^z$. The result is $\frac{2\pi i}{2!} f''(0) = \pi i$.
\end{itemize}

\section{Taylor Series and Laurent Series}

\subsection{Taylor Series}
\textbf{Real case:} Let $f: \mathbb{R} \rightarrow \mathbb{R}$ be $N+1$ times differentiable. For any $x_0 \in \mathbb{R}$, the Taylor polynomial is $T^N_{x_0} f(x)$. If $f \in \mathcal{C}^{N+1}$, there exists $\theta \in (0,1)$ such that:
\begin{align*}
    f(x) = \sum_{n=0}^N \frac{f^{(n)}(x_0)}{n!}(x-x_0)^n + \frac{f^{(N+1)}(\theta x_0 + (1-\theta)x)}{(N+1)!}(x-x_0)^{N+1}
\end{align*}

\textbf{Complex case:} Let $f: D \rightarrow \mathbb{C}$ be holomorphic. If the disk $B_R(z_0)$ is contained in $D$, then $f$ is analytic and:
\begin{align*}
    f(z) = \sum_{n=0}^\infty \frac{f^{(n)}(z_0)}{n!}(z-z_0)^n \quad \text{for all } z \in B_R(z_0)
\end{align*}
The largest such $R$ is the \textbf{Radius of Convergence}.

\subsection{Laurent Series}
\textbf{Theorem:} Let $f$ be holomorphic on an annulus (punctured disk) $0 < |z-z_0| < R$. Then:
\begin{align*}
    f(z) = \sum_{n = -\infty}^\infty c_n(z-z_0)^n = \dots + \frac{c_{-2}}{(z-z_0)^2} + \frac{c_{-1}}{z-z_0} + c_0 + c_1(z-z_0) + \dots
\end{align*}
where the coefficients are given by:
\begin{align*}
    c_n = \frac{1}{2 \pi i}\int_\gamma \frac{f(w)}{(w-z_0)^{n+1}}dw
\end{align*}

\textbf{Definition:} The series is split into two parts:
\begin{itemize}
    \item \textbf{Singular part:} Terms with negative powers ($n < 0$).
    \item \textbf{Regular part:} Terms with non-negative powers ($n \geq 0$).
\end{itemize}

\section{Classification of Singularities and Residues}

\subsection{Types of Isolated Singularities}
By looking at the singular part of the Laurent series of $f$ at $z_0$:
\begin{enumerate}
    \item \textbf{Removable Singularity:} No negative powers ($c_n = 0$ for all $n < 0$).
    \item \textbf{Pole of order $k$:} The lowest power is $c_{-k} \neq 0$ (finitely many negative terms).
    \item \textbf{Essential Singularity:} Infinitely many negative terms (e.g., $e^{1/z}$ at $z=0$).
\end{enumerate}

